

We next develop a structural model that rationalizes the spatial propagation patterns documented in the data and allows us to quantify the role of production networks in shaping the incidence of sectoral shocks. The model has a Ricardian structure that is used to recover firm-to-firm linkages not directly observed in the data. However, since the extensive margin pattern we uncover relies on regions where no suppliers comes from, we depart from the textbook Eaton-Kortum by limiting the number of varieties per sector to a fixed value. This granular Ricardian model allows us to endogeneise the region hosting no suppliers. The full derivation of the model is presented in Appendix~\ref{ap:structural}.

%==============================================================================
\subsection{Model}\label{subsec:model}
%==============================================================================



\paragraph{Environment.}
The economy consists of $R+1$ regions, where the $(R+1)^{\text{th}}$ region represents the rest of the world. \isa{we never mention how we discipline the share of input purchases sourced from the RoW} \swann{I think we do in the Appendix section called "Notes on foreign parameters". I think we can specify that we allow for imports but since we don't have the balance sheet information for those suppliers, we don't model firm-level linkages for foreign suppliers.} Firms operate either in upstream sectors $s = 1,\ldots,S$ or in the downstream industry $d$ (motor vehicles or aerospace). Each domestic region hosts at most one representative downstream firm. Firm locations and sectoral assignments are taken as exogenous. We allow for imports, but in the absence of balance-sheet information on foreign suppliers we do not model firm-level linkages abroad; the foreign share of sourcing is calibrated (Appendix~\ref{app:model-foreign}).\\ 

Domestic regions are additionally grouped into attraction areas. An attraction area $a$ is the set of upstream regions sharing the same nearest downstream producer, and we write $a(r')$ for the area of region $r'$ and $\mathcal{R}_{a}$ for its members. Because the model is estimated separately for the motor vehicles and aerospace industries, we omit the downstream-industry index in the remainder of the exposition. Unless otherwise stated, all variables and parameters should therefore be understood as industry-specific.


\paragraph{Final demand.}
Downstream firms compete under monopolistic competition and face an isoelastic demand function:
$$q_{r} = \left(\frac{p_{r}}{P}\right)^{\varepsilon-1}\frac{E}{P}\,\delta_{r}^D,$$
where $\delta_{r}^D$ is a region-specific demand shifter, $E/P$ denotes aggregate real expenditure and $p_{r}/P$ is the relative price of the variety produced in region $r$.
Demand shifters are later used to simulate downstream shocks and trace their propagation through production networks. The market for final goods is assumed to be perfectly integrated.

\paragraph{Production.}
Downstream firms produce using labor and intermediate inputs according to a constant-returns-to-scale production technology:
\begin{equation}
  y_r = A_r g(\ell_r,\{q_{rs}\}_{s=1}^{S})  
\end{equation}
where $A_r$ denotes total factor productivity, $\ell_r$ is labor input and $q_{rs}$ is the quantity of the input bundle from sector $s$ used in production. As discussed in Appendix~\ref{apsec:structural}, the structural estimation assumes a nested CES specification, where $\Omega^L$ scales the labor share, $\Omega^s$ denotes the expenditure share on inputs from sector $s$ and each $q_{rs}$ is itself a CES aggregate over a of a fix number $N_s$ of varieties.

\paragraph{Upstream technology and sourcing.}
Each upstream sector $s$ consists of $N_{s}$ horizontally
differentiated varieties, indexed $\rho=1,\dots,N_{s}$. For each variety, there is an infinite number of firms producing it with a constant-returns-to-scale technology that relies solely on labor. Productivity draws follow a sector- and region-specific Fréchet distribution:
$$\Pr\left(z_{rs}(j) \leq z\right) = e^{-T_{rs}z^{-\theta}}, \quad T_{rs} > 0.$$
The parameter $T_{rs}$ governs average productivity, while $\theta$ controls the dispersion of productivity draws. Shipping intermediate goods from region $r'$ to region $r$ involves iceberg trade costs, denoted $\tau_{r'rs}$ which are parameterized as a function of geographic distance.

In the baseline specification, comparative advantage is a property of the attraction area rather than of the individual sector x region. \begin{equation}
T_{r's}=T_{a(r')s}\qquad\text{for every }r'\in\mathcal{R}_{a}.
\label{eq:T-area}
\end{equation}

The model yields two important quantities that will use to discipline the parameters in the estimation. First are the bilateral sourcing shares. Following \cite{EatonKortum2002}, downstream firms source each variety from the lowest-cost supplier, in France or abroad. The probability that a downstream firm located in region $r$ sources an input from sector $s$ in region $r'$ is given by:
\begin{equation}
\gamma_{r'rs} = \frac{T_{r's}(w_{r's}\tau_{r'rs})^{-\theta}}
                     {\sum_{k=1}^{R+1} T_{ks}(w_{ks}\tau_{krs})^{-\theta}}
\label{eq:gamma_rrs}
\end{equation}
where $w_{r's}$ denotes the unit production cost in region-sector pair $(r',s)$. Because trade costs increase with geographic distance, sourcing probabilities decline with spatial separation, generating localized production networks. 

We then aggregate bilateral sourcing shares at the regional level:

\begin{equation}\label{eq:gamma_rs}
    \gamma_{r's}=\sum_r \frac{X_{rs}}{ \sum_r X_{rs}} \gamma_{r'rs}
\end{equation}
where $X_{rs}$ denotes the value of intermediate inputs from sector $s$ shipped to downstream firms located in region $r$. This quantity will be used to pinned down comparative advantages. 


Second is the extensive margin probability. Our surveys record whether an upstream firm supplies the downstream industry, not which buyer it serves. The relevant object is therefore not the bilateral sourcing probability \eqref{eq:gamma_rrs} but the probability that a region wins a given variety somewhere in the industry --- the union over downstream buyers,

\begin{equation}
p_{r's}\;\equiv\;\Pr\Big(\bigcup_{r=1}^{R}\big\{\text{$r'$ wins variety $\rho$ for buyer $r$}\big\}\Big).
\label{eq:p}
\end{equation}




This extensive margin probability falls with the distance and raises with comparative advantages. For a given variety the champion cost is common across buyers and only $\tau_{r'rs}$ differs, so win events are
positively dependent across buyers and $p_{r's}>1-\prod_{r}(1-\gamma_{r'rs})$: a firm productive enough to win one market is more likely to win the others. For a large number of downstream regions, there is no tractable closed form for $p_{r's}$, and we evaluate it by simulation. This is why the model is estimated by simulated rather than analytical method of moments.\footnote{Writing $p_{r's}$ in closed form requires inclusion--exclusion over subsets of downstream buyers,
$p_{r's}=\sum_{\emptyset\neq\mathcal{S}\subseteq\{1,\dots,R\}}(-1)^{|\mathcal{S}|+1} \Pr(\cap_{r\in\mathcal{S}}\{\text{$r'$ wins for }r\})$, a sum with $2^{R}-1$ terms. With the number of downstream regions we observe this is prohibitive, and dominance pruning of the subsets --- discarding buyers whose win event contains another's --- reduces the count but not reliably enough to be usable inside an optimizer. }

With $N_{s}$ varieties the win events are independent across $\rho$, so the industry-wide supplier count is binomial,

\begin{equation}
K_{r's}\sim\mathrm{Bin}\big(N_{s},p_{r's}\big), \qquad \Pr\big(K_{r's}=0\big)=(1-p_{r's})^{N_{s}}>0 .
\label{eq:binomial}
\end{equation}

The zero-supplier probability is strictly positive and endogenous: it falls with comparative advantage, rises with distance through, and falls with the variety count: when $N_s \rightarrow 0$ we fall back to the Eaton-Kortum model where a downstream firm sources from all regions. For each sector, we define $G_s(K)$ the share of regions with at most $K$ suppliers.




%==============================================================================
\subsection{Estimation and fit}\label{subsec:estimation_fit}
%==============================================================================

We estimate the model using Simulated Method of Moments. Unlike the standard Ricardian trade literature, which relies on bilateral input–output flows, such flows are not observed in our setting. Instead, we exploit survey information on supplier relationships together with VAT data on firm sales to discipline the spatial structure of sourcing and demand. These moments provide direct information on the geography of production networks and allow us to identify the key parameters governing trade costs and regional comparative advantage in the \cite{EatonKortum2002} framework.

We parametrize the trade cost using the distance between two regions: $\tau_{r'rs}=d_{r'r}^{\alpha}$; and obtain the model's parameters:
\begin{equation}
\big\{\Omega_{L},\;\{\Omega_{s}\},\;\alpha,\;\{A_{r}\},
\{T_{as}\},\;\{N_{s}\}\big\},
\end{equation}

Beyond these structural parameters, the model includes several demand- and supply-side elasticities, which we calibrate as described in Appendix~\ref{app:model-calibration}. We also provide there additional details on the optimization procedure and on the calculation of standard errors. 


\paragraph{Targeted moments.} The model is matched to six families of moments. The downstream industry's labor share and the sectoral composition of intermediate expenditures are constructed from input-output tables and identify $\Omega_{L}$ and $\{\Omega_{s}\}_{s=1}^{S}$. The regional distribution of downstream sales, $\pi_{r}$, is recovered from the VAT data and identifies the downstream productivity parameters $\{A_{r}\}$, up to a multiplicative constant and therefore normalized relative to the largest region.

The elasticity of trade costs to distance is discipline by the extensive margin (Equation~\eqref{eq:reduced_form_spatialIO}). Contrary to this section, we don't estimate the profile using a linear probability model. We use instead the prediction of the model in the limit of a single downstream buyer $a$. For a firm $z$ located in $(r',s)$ its probability to win variety $\rho$'s market is $\rho_{r's}(z) = e^{-\Phi_{as}(w_{r's}\tau_{r'rs})^{\theta}z^{-\theta}}$ which implies that $\Pr(\text{no supply} \mid z) = 1 - \exp(-\exp(\eta))$ where $\eta = \ln \Phi_{as} + \theta \ln w_{r's} + \theta \alpha \ln d_{r'a} - \theta \ln z$. Consequently, we are going to estimate both in the data and the model the following complementary log log design:

\begin{equation}\label{eq:firm-cloglog}
\Pr(\text{no supply} \mid z) = 1 - \exp(-\exp(FE_{a(r')s} + \theta \alpha \ln d_{r'a(r')} - \theta \ln z))  
\end{equation}



where $z$ is the productivity of the firm, approximated in the data by its sales\footnote{More precisely, we run two specifications in the data. One where we regress on the distance as a continuous variable and use the estimated coefficient as a starting point for $\alpha$ in the estimation. In the second, the distance is binarized as in the linear probability model to build the targeted moments.}. 

Contrary to the first set of parameters, comparative advantages and number of varieties are not free parameters. They are calibrated inside the model. For $\{T_{as}\}$, the area-aggregated sourcing shares are as numerous as the comparative-advantage parameters, and for given trade costs and expenditures the mapping from $T$ to those shares is invertible: there is a unique vector of comparative advantages, up to the per-sector reference normalization, that reproduces the observed shares exactly. We compute it by a matrix-scaling iteration alternated with the general-equilibrium solve (Appendix~\ref{app:model-estimation}). For $\{N_{s}\}$, the count moment $\bar G_{s}(0)$ is decreasing in the variety count and, by the argument below, does not require re-simulation, so the variety count that best matches it is found by a monotone integer search over the bounds \eqref{eq:app-bounds}. 


\paragraph{Identification.} The channels are close to orthogonal by construction. $\alpha$ is read from a within-area gradient that comparative advantage cannot touch, because comparative advantage is constant within the area. $T_{as}$ is read from across-area value shares. $N_{s}$ is read from the level of the empty-region share, which the fixed effect strips out of the $\alpha$ regression. To validate our moments selection, we perform a sensitivity analysis in Appendix~\ref{app:sensitivity_analysis}. As we have many parameters, we only provide the Jacobian of the model, evaluated at the calibrated parameter set.  

\begin{figure}[H]
    \centering
\caption{Empirical vs. simulated moments\label{fig:emp_sim_gamma_coef_G0}}
    \begin{tabular}{cc}
    Motor vehicles &Aerospace\\
    \multicolumn{2}{c}{Regional sourcing shares ($\gamma_{as}$)}\\
        \includegraphics[width = 0.5\linewidth]{figures/emp_sim_gamma_aa_auto_mu1.pdf}&\includegraphics[width = 0.5\linewidth]{figures/emp_sim_gamma_aa_aero_mu1.pdf}\\ \vspace{1em}\\
\multicolumn{2}{c}{Extensive margin coefficients}\\
 \includegraphics[width = 0.5\linewidth]{figures/emp_sim_reg_coef_auto_mu1.pdf} &  \includegraphics[width = 0.5\linewidth]{figures/emp_sim_reg_coef_aero_mu1.pdf}  \\   \vspace{1em}\\  
 \multicolumn{2}{c}{Share of commuting zones without supplier ($G_s(0)$)}\\
 \includegraphics[width = 0.5\linewidth]{figures/emp_sim_G0_auto_mu1.pdf} &  
 \includegraphics[width = 0.5\linewidth]{figures/emp_sim_G0_aero_mu1.pdf}  \\  
    \end{tabular}
    \caption*{\footnotesize \emph{Notes}: The figure displays scatter plots of empirical versus simulated regional sourcing shares ($\gamma_{r's}$, top panels) and the extensive margin coefficients estimated using complementary log log model (middle panels) and share of commuting zone without supplier (bottom panel) for the motor vehicles industry (left panels) and the aerospace industry (right panels). The black line denotes the 45-degree line. The statistics reported in each top panel correspond to the slope coefficient from a linear regression of simulated moments on empirical moments. 
    In the top panel, regional sourcing shares for the reference attraction areas are in red. For the middle and bottom panel, standard errors for the empirical points are computed using bootstrap. Standard errors for the simulated points are recovered from the delta method (see Appendix~\ref{app:model-standard_errors}). 
    }
\end{figure}

Figure~\ref{fig:emp_sim_gamma_coef_G0},  Table~\ref{tab:moments_comparison_combined} and Figure~\ref{fig:emp_sim_pi} in the Appendix compare empirical and simulated moments for both industries. The model reproduces closely the aggregate labor share and sectoral expenditure shares (Table~\ref{tab:moments_comparison_combined}), including the strong within-sector sourcing observed in aerospace (62\%) and the prominent role of automotive equipment in motor vehicles (24\%). Similarly, the model tracks closely the distribution of downstream sales (Figure~\ref{fig:emp_sim_pi}) for both industries.  \swann{En réalité le fit de la distribution des ventes n'est pas top pour les petites downstream. Je laisse ça comme ça comme ça risque de changer. }

In the top panel of Figure~\ref{fig:emp_sim_gamma_coef_G0} we report the comparison for the regional sourcing shares $\gamma_{as}$ at the attraction area level. Simulated moments align closely with the 45-degree line in both industries. In the middle panel, the extensive margin coefficients implied by the model are most of the time not statistically different from the one observed in the data. Finally, the bottom panel shows that the empirical and simulated distribution of the count moment align closely. The number of varieties per sector range from 2 to 23 in the Motor Vehicles industry and from 3 to 14 in the Aerospace industry. \\

This Ricardian model has a novel property, the fixed number of varieties per sector, that we disciplined using the count moment $G_s(0)$. To provide additional evidence of the quality of our estimation, we provide in Appendix Figure~\ref{app:untargeted_count} the fit of higher order count moment $G_s(K)$ for $K \in \{1,2,3\}$. Compared to the data, the model is good at reproducing the share of region with 2 or 3 suppliers but overstate those with exactly one supplier. \\
\begin{comment}
    
\paragraph{Analyzing the incidence of shocks.} In the remainder of this paper, we use the estimated model to analyze how shocks originating from downstream regions propagate through the production network. Because the model features a finite number of varieties, each network realization depends on the specific draw of Ricardian productivity shocks. Consequently, our analysis distinguishes between systematic sourcing patterns and the variation arising from this granularity. We consider both the expected network, which captures the systematic patterns implied by the model, and the distribution of actual network realizations. For linear statistics, such as average sourcing distance, the expected statistic coincides with the statistic computed on the expected network. For nonlinear statistics, however, the results reflect both systematic patterns and the dispersion generated by finite granularity.\footnote{The expectation $\mathbb{E}_{\omega}$ is taken over independent realizations $\omega$ of the Ricardian productivity draws. While linearity implies $\mathbb{E}_{\omega}[f(N(\omega))]=f(\mathbb{E}_{\omega}[N(\omega)])$, this equality generally fails for nonlinear statistics (e.g., variances, concentration indices, or effective numbers of suppliers) because they depend on higher moments of the sourcing shares.}

\paragraph{Trade costs, comparative advantage, and spatial sorting.}
The estimated model attributes the geography of sourcing to two forces: comparative advantage across upstream regions and the cost of trading across distance. Disentangling the two from sourcing patterns alone is difficult because a region may be a large supplier either because it is particularly productive or because it is close to demand. The two forces are observationally similar unless the trade-cost elasticity $\alpha\theta$ is identified separately. This is precisely what the estimation delivers, allowing us to assess both the relative strength of the two forces and their spatial alignment.

A first comparison concerns the dispersion of comparative advantage. Averaged across sectors, the within-sector standard deviation of $\log \hat T_{as}$ is nearly identical across the two industries: 0.91 in motor vehicles and 0.96 in aerospace. The industries therefore exhibit similar heterogeneity in comparative advantage. They differ instead in how strongly sourcing responds to distance: the estimated trade-cost elasticity is $\alpha\hat\theta=0.23$ in motor vehicles and $0.46$ in aerospace. Thus, conditional on the same distribution of distances, geography generates substantially more dispersion in sourcing attractiveness in aerospace. A comparative-advantage gap $\Delta$ is offset by a distance ratio $\exp(\Delta/\alpha\theta)$, so a smaller $\alpha\theta$ requires a larger distance difference to offset a given productivity advantage. In this sense, an identical comparative-advantage gap has a stronger effect on sourcing in motor vehicles than in aerospace.

The previous comparison considers the two forces separately. To assess how they jointly shape the allocation of sourcing across competing locations, consider the sourcing index $\log\psi_{r'rs} = \log T_{a(r')s}-\alpha\theta\log d_{r'r}.$ For each sector $s$ and buyer $r$, we compute the variance of $\log\psi_{r'rs}$ across the supplier cells $l$ competing for that buyer, namely:

$$
\operatorname{Var}_{r'}(\log\psi_{r'rs}) = \operatorname{Var}_{r'}(\log T_{a(r')s}) + (\alpha\theta)^2\operatorname{Var}_{r'}(\log d_{r'r}) - 2\alpha\theta\operatorname{Cov}_{r'}(\log T_{a(r')s},\log d_{r'r}).
$$

The first term measures dispersion in comparative advantage across potential suppliers, the second measures the dispersion generated by distance alone, and the third captures their spatial alignment. In particular, when high-$T$ suppliers tend to be closer to the buyer, $\operatorname{Cov}(\log T,\log d)<0$, so the covariance term is positive: comparative advantage and proximity reinforce one another rather than offsetting each other. For each sector, we aggregate these buyer-specific variances using the buyer's purchase weights. The percentages in Table~\ref{tab:variance_decomposition} therefore describe the contribution of each component to the estimated cross-supplier dispersion in $\log\psi$, after accounting for the sourcing weights of buyers and the distribution of sectors.

\begin{table}[htbp]
\centering
\caption{Variance decomposition of the sourcing index}
\label{tab:variance_decomposition}
\begin{tabular}{lcc}
\toprule
& Motor vehicles & Aerospace \\
\midrule
Comparative advantage
& 87\% & 86\% \\
Distance
& 2\% & 7\% \\
Co-location
& 11\% & 9\% \\
\midrule
\bottomrule
\end{tabular}
\caption*{\emph{Notes:} The decomposition is computed across competing supplier cells for each buyer and sector. Buyer-specific variances are aggregated using purchase weights and then averaged across sectors. The comparative-advantage component is $\operatorname{Var}(\log \hat T)$, the distance component is $(\alpha\hat\theta)^2\operatorname{Var}(\log d)$, and the co-location component is $-2\alpha\hat\theta\operatorname{Cov}(\log \hat T,\log d)$. Percentages are normalized by the corresponding variance of $\log\hat\psi$.}
\end{table}

The decomposition yields three conclusions. First, comparative advantage is by far the dominant source of dispersion in sourcing attractiveness in both industries: it accounts for 87\% of the variance in motor vehicles and 86\% in aerospace. The industries do not differ in how unequal supplier productivity is, in its dispersion or in its weight in the allocation.

Second, they separate on the \emph{alignment} term rather than on distance. The pure distance component accounts for 2\% of the variance in motor vehicles and 7\% in aerospace --- the ordering $\alpha\hat\theta$ requires --- but the co-location component contributes 11\% in motor vehicles, more than five times the contribution of distance alone, against 9\% in aerospace, only modestly above the 7\% distance already carries. Scale-free, the same statement is $\operatorname{Corr}_l(\log\hat T,\log d)=-0.47$ at the median motor-vehicle sector against $-0.17$ in aerospace. Productive suppliers sit close to buyers in motor vehicles; in aerospace comparative advantage and proximity are much closer to being independent sources of sourcing variation.

Third --- and this is what a median cannot show --- the two industries differ in the \emph{uniformity} of that alignment, not only in its strength. Figure~\ref{fig:alignment_correlation} reports the distribution of $\rho_{rs}=\operatorname{Corr}_l(\log\hat T_{a(l)s},\log d_{lr})$ across all (sector, buyer) pairs. In motor vehicles the sourcing-weighted covariance is negative in every sector, between $-6$ and $-116$~km per log point with a median of $-54$, and the alignment is one-sided: only 19\% of (sector, buyer) pairs carry the opposite sign, so hardly any buyer faces a geography in which comparative advantage and proximity pull apart. In aerospace the median is $-14$~km and the sign is not even common across sectors: other transport equipment, which alone carries 56\% of the industry's modelled intermediate purchases, has a \emph{positive} covariance of $+44$~km. A sector median of $-0.17$ there is an average over buyers pulled in opposite directions, not a mild version of the motor-vehicle case; the two are different economies, and they will predict different counterfactuals.


Two consequences of the estimated $\alpha\hat\theta$ are worth stating before the counterfactuals, because they govern how the geography of sourcing can differ across the two industries at all. The distance factor $d_{lr}^{-\alpha\theta}$ is the \emph{only} term in the sourcing probability that varies with the identity of the buyer: comparative advantage is a property of the supplier's area alone. A low $\alpha\hat\theta$ therefore makes sourcing less buyer-specific --- over the range of French bilateral distances, the most remote supplier is only $2.4$ times less attractive than the closest one in motor vehicles, against $5.7$ times in aerospace --- so buyers in motor vehicles draw from close to a common national distribution while aerospace buyers tilt sharply towards their own neighbourhood. And a negative $\operatorname{Cov}(\log T,\log d)$ says productive suppliers sits near demand, but it does not say whether they cluster in one or several regions. Section~\ref{subsec:counterfactuals} settles which of the two the estimates describe, by propagating an actual shock through the estimated allocation.

\end{comment}
%==============================================================================
\section{Quantitative implications for spatial comovements}\label{subsec:counterfactuals}
%==============================================================================

In the remainder of this paper, we use the estimated model to analyze how shocks originating from downstream regions propagate through the production network. Especially, we try to understand the role of comparative advantages, trade cost and granularity in shaping this propagation. \\

Because the model features a finite number of varieties, each network realization depends on the specific draw of Ricardian productivity shocks. Consequently, our analysis distinguishes between systematic sourcing patterns and the variation arising from this granularity. We consider both the expected network, which captures the systematic patterns implied by the model, and the distribution of actual network realizations. For linear statistics, such as average sourcing distance, the expected statistic coincides with the statistic computed on the expected network. For nonlinear statistics, however, the results reflect both systematic patterns and the dispersion generated by finite granularity.\footnote{The expectation $\mathbb{E}_{\omega}$ is taken over independent realizations $\omega$ of the Ricardian productivity draws. While linearity implies $\mathbb{E}_{\omega}[f(N(\omega))]=f(\mathbb{E}_{\omega}[N(\omega)])$, this equality generally fails for nonlinear statistics (e.g., variances, concentration indices, or effective numbers of suppliers) because they depend on higher moments of the sourcing shares.}

The propagation of the shock is studied through four angles: the length of the diffusion, its geographical concentration across buyers and its local incidence. The results reveal two distinct dimensions of spatial shock propagation. First, comparative advantage determines where production linkages are concentrated, while trade costs determine how far a shock travels from its origin. These forces generate markedly different network architectures across industries. In motor vehicles, comparative advantage and proximity reinforce each other, producing a dispersed network of suppliers that is broadly shared across buyers: shocks propagate widely, but in similar ways across regions. In aerospace, by contrast, comparative advantage concentrates sourcing in a few common hubs, while distance determines how much of the shock remains local; as a result, shocks originating outside the main hubs are transmitted over long distances toward them, whereas shocks originating in the hubs are largely absorbed locally. The same distinction shapes local incidence: comparative advantage reinforces local propagation in motor vehicles but pulls activity away from the originating region toward the aerospace hubs. Thus, production networks affect not only the size of the response to a demand shock, but also where that response is concentrated and whether it remains local or is transmitted across space.




%The geographic incidence is first studied through the average sourcing distance which represent the average distance a shock hitting a buyer will propagate. We quantify how this distance moves when we equalize comparative advantages across space or shut down trade costs. Such is limited when one want to understand the geography of the propagation of the shock. How strongly are sourcing portfolio across supplier geographically concentrated? We answer this question comparing the geography of the portfolio to an homogeneous distribution of suppliers. Finally, shock propagation might not be a linear function of distance and we investigate the local propagation of those. 

%Armed with the estimated model, we can now study the diffusion of shocks from downstream producers to upstream suppliers. Given a localized demand shock, we quantify (i) the amplification multiplier associated with production linkages and (ii) its spatial allocation. While the former arises in any model with input-output linkages and is primarily determined by the production-function parameters ($\Omega^L, {\Omega^s}$), the latter is more specific to our framework with endogenous production linkages. In our model, the spatial diffusion of shocks is governed by the interaction between regional comparative advantage (${T_{rs}}$) and trade costs (${\tau_{r'rs}}$).

\begin{comment}
We organize this section around two exercises: 

\begin{enumerate}
    \item The spatial diffusion of a €1 downstream shocks as a function of its origin. We first quantify the amplification multiplier and spatial heterogeneity of diffusion of the shock within a 100\km~ radius centered on the shocked region. Then we embedded the model prediction in order to analyze recent shocks affecting both industries.  For aerospace, the 2019--2025 Military Programming Law; for motor vehicles, the 2024 EV eco-score subsidy reform.
    \item The consequence of a shock on trade costs or comparative advantage in the reallocation of the production network and the comovement between regions. 
\end{enumerate}


\swann{A réécrire. }
\end{comment}
\subsection{Spatial diffusion of downstream shocks.}


\paragraph{Amplification.} We first quantify how much production linkages amplify a demand shock. The amplification coefficient $D_{r}$ is the aggregate increase in sales generated by one additional euro of final demand in a given downstream region. In our framework it has a closed form,
$$D_{r}=1+(1-\Omega^{L})\left(P_{r}/c_{r}\right)^{1-\lambda},$$
that is, the euro itself plus the share of the downstream unit cost spent on domestic suppliers. One euro of final demand generates 85 cents of additional upstream sales in motor vehicles and 87 cents in aerospace, with almost no spatial variation. The homogeneity is mechanical. Technical input shares are common within a downstream industry, and neither comparative advantage nor trade costs enter $D_{r}$; the only regional term is local access to cheap inputs, through $P_{r}/c_{r}$. \isa{Would it be possible to compare with all sectors (eg in appendix), using solely the IO coefficients? if it is not too difficult, that would help say something about the scale: Is 1.5 high or not? How big is it compared to the aggregate effect, which can be recovered from the Leontief inverse? The simple IO case is described in details in Acemoglu et al (2016): \url{https://www.journals.uchicago.edu/doi/epdf/10.1086/685961} Note that they also have a section on the local diffusion of shocks, which we should cite} \isa{is it true? My intuition would be that a relatively difficult access to low price suppliers generates i) lower input purchases (substitution to labor) and ii) more input purchases to non-French firms. Both channels should reduce the multiplier. However, I find it surprising that we find among the lowest amplification multiplier terms a region like Toulouse which should in principle get good access to local inputs. Is the reason to be found in low elasticities of substitution, that generate a positive impact of relatively expansive input purchases? }

\subsubsection{Length of the diffusion.}
How far does a shock hitting a buyer in region $r$ travel? We measure it by the average sourcing distance, the distance travelled by the average euro of upstream purchases:
\begin{equation}
\bar d_{r}=\frac{\sum_{r'} X_{r'r}\,d_{r'r}}{\sum_{r'} X_{r'r}}.
\label{eq:sourcing_distance}
\end{equation}
In levels, the two industries look alike: the average buyer sources at 348~km in motor vehicles and 355~km in aerospace. Levels, however, mix the allocation of sourcing with the shape of the country: a buyer in Brittany faces longer distances than a buyer in the centre of France under any allocation, and aerospace buyers are more spatially dispersed than motor-vehicle buyers. We therefore divide $\bar d_r$ by the distance the same euro would travel if sourcing were unselective, that is, if neither comparative advantage nor trade costs discriminated among potential suppliers.\footnote{Let $\mathcal{L}_s$ denote the set of modelled domestic regions supplying sector $s$ and $X_r=\sum_s X_{rs}$. The uniform benchmark sets $\alpha=0$ \emph{and} equalises $\hat T_{as}$, so that each region in $\mathcal{L}_s$ receives the same share of the sector's spending: $\bar d_r^{\,U}=\sum_s \frac{X_{rs}}{X_r}\,\frac{1}{|\mathcal{L}_s|}\sum_{r'\in\mathcal{L}_s} d_{r'r}$. 
It is a geometric reference, not a counterfactual economy. The normalised statistic $\bar d_r/\bar d_r^{\,U}$ ranges over $[0.68,\,1.08]$ in motor vehicles and $[0.64,\,0.98]$ in aerospace, against raw ranges of 257--407~km and 253--590~km.} The normalised statistic separates the industries: the median motor-vehicle buyer sources at $0.78$ of the unselective distance, the median aerospace buyer at $0.91$. Selection pulls motor-vehicle sourcing substantially closer to the buyer, and aerospace sourcing barely at all.

Since the average sourcing distance is a linear statistic of the network, only comparative advantage and trade costs influence it. We switch them off one at a time: equalising comparative advantage across areas (\emph{Distance only}) or setting $\alpha=0$ (\emph{Comparative advantage only}). In both counterfactuals, sectoral spending $X_{rs}$ is held fixed and only its allocation across regions changes.\footnote{Wages are normalised to one, so $T_{as}$ should be read as $w^{-\theta}T_{as}$. Foreign sourcing is not reallocated.} Both forces can be dialled continuously. Writing the sourcing probability as

\begin{equation*}
\tilde\gamma_{r'rs}(t)=\frac{T_{a(r')s}^{\,t}\,d_{r'r}^{-\alpha\theta}}{\Phi_{rs}(t)},
\qquad
\Phi_{rs}(t)=\sum_{k\in\mathcal{L}_s}T_{a(k)s}^{\,t}\,d_{kr}^{-\alpha\theta},
\end{equation*} 
the estimated allocation corresponds to $t=1$ and \emph{Distance only} to $t=0$. The average sourcing distance decomposes across sectors as $\bar d_r=\sum_s (X_{rs}/X_r)\,\bar d_{rs}$, with $\bar d_{rs}=\sum_{r'\in\mathcal{L}_s}\tilde\gamma_{r'rs}\,d_{r'r}$. Because the sourcing probabilities form an exponential family, each force moves $\bar d_{rs}$ at a rate given by a single covariance, taken across the buyer's potential suppliers and weighted by their sourcing probabilities:

\begin{align}
\frac{\partial \bar d_{rs}}{\partial t}
&=\operatorname{Cov}_{\tilde\gamma_{\cdot rs}}\!\big(\log T_{a(r')s},\,d_{r'r}\big),
\label{eq:alignment}\\
\frac{\partial \bar d_{rs}}{\partial(\alpha\theta)}
&=-\operatorname{Cov}_{\tilde\gamma_{\cdot rs}}\!\big(\log d_{r'r},\,d_{r'r}\big)\;<\;0.
\label{eq:distance_rate}
\end{align}
The full effect of each counterfactual integrates these rates along the path; for comparative advantage, $\bar d_r(0)-\bar d_r(1) =-\int_0^1\sum_s (X_{rs}/X_r)\operatorname{Cov}_{\tilde\gamma_{\cdot rs}(t)} (\log T_{a(r')s},d_{r'r})\,dt$.

The two expressions have very different implications. Since $d$ and $\log d$ always move together, the covariance in \eqref{eq:distance_rate} is positive for every buyer and sector: trade costs \emph{always} shorten the diffusion, and removing them lengthens it everywhere. Comparative advantage has no such sign restriction. By \eqref{eq:alignment}, $\hat T$ affects the sourcing distance only through its alignment with the buyer's geography. Where productive suppliers sit close to the buyer ($\operatorname{Cov}<0$), comparative advantage shortens the diffusion, and equalising it sends the euro further. Where they sit far away ($\operatorname{Cov}>0$), comparative advantage lengthens it. The dispersion of $\hat T$ matters only through this covariance: a highly unequal $\hat T$ that is unrelated to the buyer's geography reallocates sourcing across regions without changing how far the euro travels.

Figure~\ref{fig:alignment_covariance} reports the distribution of the covariance in \eqref{eq:alignment} across (sector, buyer) pairs, in kilometres per log point of comparative advantage. Trade costs shorten the diffusion in both industries, by 23~km (6.6\%) in motor vehicles and 34~km (9.6\%) in aerospace, and they do so for \emph{every} buyer in both --- all 22 motor-vehicle and all 18 aerospace buyers source further once $\alpha$ is set to zero, as \eqref{eq:distance_rate} requires.

Comparative advantage is where the two industries part, and it parts them on the \emph{dispersion} of that rate rather than on its centre. Both distributions sit on the negative side --- a median of $-48$~km per log point in motor vehicles and $-11$ in aerospace --- but they differ in how much of the mass lies beyond zero, and in how far it lies. Only 19\% of motor-vehicle (sector, buyer) pairs carry a positive covariance, against 42\% of aerospace ones, and the aerospace support is both wider and nearly symmetric: a tenth-to-ninetieth-percentile range of $[-92,\,+96]$~km against $[-126,\,+18]$ in motor vehicles. In motor vehicles almost every pair pulls sourcing towards the buyer. In aerospace two fifths pull it away, and by as much as the rest pull it in.

Figure~\ref{fig:counterfactual_distance_region} shows what those rates produce once they are aggregated with each buyer's spending and integrated along the path, reporting each buyer's sourcing distance as a percentage deviation from the estimated allocation. In motor vehicles a one-sided distribution of rates gives a one-sided displacement: 20 of the 22 buyers source further when $\hat T$ is equalised, and the average sourcing distance rises by 65~km, or 19\%, close to three times the 6.6\% that trade costs are worth. In aerospace the two directions very nearly cancel in the aggregate --- the average moves by 2~km, under one per cent --- while exactly half the buyers source further and half source closer. What cancels is large, not small: the average buyer's sourcing distance moves by 8.8\% in absolute value against a net of 1.3\%, and the individual deviations run from $+30\%$ to $-14\%$. Buyers in Toulouse lose the advantage of their local suppliers, and their sourcing distance rises by 30\%. Buyers in \^Ile-de-France, whose purchases were drawn towards distant hubs such as Toulouse, source more locally, and their sourcing distance falls by 11\%. Trade costs admit no such split: with $\alpha$ set to zero every one of the 18 buyers sources further, by 10\% on average and by at most 25\%.

Comparative advantage thus shapes the length of the diffusion in both industries, but through opposite geographies. In motor vehicles, it reinforces proximity. In aerospace, it pulls sourcing towards a few hubs that are close to some buyers and far from others.
\begin{figure}[H]
    \centering
\caption{The rate at which comparative advantage moves the sourcing distance\label{fig:alignment_covariance}}
    \includegraphics[width = 0.65\linewidth]{figures/alignment_covariance_mu2.pdf}
    \caption*{\footnotesize \emph{Notes}: Distribution across (sector, downstream buyer) pairs of $\operatorname{Cov}_{\tilde\gamma_{\cdot rs}}(\log \hat T_{a(l)s},\,d_{lr})$, in kilometres per log point of comparative advantage: by \eqref{eq:alignment} this is the rate at which dialling $\hat T$ down moves that buyer's average sourcing distance in that sector, and its path integral is the \emph{Distance only} counterfactual. The moments are taken across the supplier cells $l$ competing for that buyer and weighted by the sourcing probabilities, which is the measure the identity holds under. A kernel density is drawn over the histogram; it is an estimate with a bandwidth, and the outline underneath is the raw count it should be read against. Mass to the left of zero is comparative advantage sitting near demand.}
\end{figure}
 
\begin{figure}[H]
    \centering
\caption{Average sourcing share per buyer under each scenario, \label{fig:counterfactual_distance_region}}
    \begin{tabular}{cc}
    Motor vehicles&Aerospace\\
        \includegraphics[width = 0.5\linewidth]{figures/counterfactual_distance_region_auto_mu2.pdf} &
        \includegraphics[width = 0.5\linewidth]{figures/counterfactual_distance_region_aero_mu2.pdf}
    \end{tabular}
    \caption*{\footnotesize \emph{Notes}: The average sourcing distance $d_{r}$ of each shocked commuting zone under each counterfactual regime, expressed as a percentage deviation from \emph{Both forces}, the estimated allocation; the level it is taken against is printed in kilometres beside each row. Percentages rather than kilometres because every $d_{r}$ lies between roughly 250 and 450~km, so on a kilometre axis the level --- common across regions and carrying no information --- crowds out the movement. Regions are in the order of Figure~\ref{fig:counterfactual_local_share}, so the two figures are read row for row.}
\end{figure}
 
\subsubsection{The geography of shock propagation.}
The average sourcing distance says how far a shock travels, not where it lands. A buyer sourcing at 350~km may be drawing on a ring of neighbours or on a single distant hub, and for the stability of the industry the two are not the same economy. What distinguishes them are two questions, and they are the same object read from its two sides. From the buyer's side: \emph{how concentrated is a buyer's portfolio of suppliers?} From the supplier's side: \emph{how concentrated is a supplier region's portfolio of customers?}. We address those questions first by analyzing the determinant of the concentration of the portfolio of buyers. Then, we investigate the overlap of those portfolio. 

\paragraph{The portfolio of a single buyer.}
Let $\Omega$ denote one realisation of the network. The portfolio of sector-$s$ suppliers of buyer $r$ is the vector of shares
\begin{equation}
\omega_{rs}(\Omega)=\left(\frac{X_{r'rs}(\Omega)}{\sum_{r'}X_{r'rs}(\Omega)}\right)_{r'},
\label{eq:incidence_vector}
\end{equation}
which is also the incidence vector of a shock to that buyer: the share of the upstream response landing in each supplier region. We are going to use the sector $\times$ buyer ($H_rs$) and buyer specific measure of portfolio concentration ($H_r$)\footnote{$H_{r}$ is the purchase-weighted average of \emph{within-sector} Herfindahls and not the Herfindahl of the pooled portfolio $\omega_{r}(\Omega)=\sum_{s}(X_{rs}/X_{r})\omega_{rs}(\Omega)$; by convexity the first is the larger of the two, and the gap is the diversification a buyer obtains from the fact that its ten input sectors have different supplier geographies. The decomposition below requires the sector-by-sector object, so that is what we report, but $1/H_{r}$ should be read as an effective number of supplier regions \emph{per sector} rather than as the number of places a buyer's whole input bill reaches. The count of places is \eqref{eq:region_reach} below.}:

$$H_{rs}(\Omega)=\sum_{r'}\omega_{r'rs}^{2}(\Omega)=\big\|\omega_{rs}(\Omega)\big\|^{2}, \quad H_{r}(\Omega)=\sum_{s}(X_{rs}/X_{r})H_{rs}(\Omega)$$
Unlike the average sourcing distance, $H$ is a quadratic statistic of the network. A linear statistic depends on the realisation only through its mean, so the expected network $\gamma_{rs}$ would summarise it entirely; a quadratic one also depends on the dispersion of the realisation around that mean. Granularity therefore enters $H$ as a term of it, alongside comparative advantage and trade costs.

That is the content of the first decomposition. Taking expectations over realisations and completing the square coordinate by coordinate,
\begin{equation}
\mathbb{E}_{\Omega}\big[H_{rs}(\Omega)\big]
=\underbrace{\big\|\gamma_{rs}\big\|^{2}}_{\text{structural}}
+\underbrace{\mathbb{E}_{\Omega}\big\|\omega_{rs}(\Omega)-\gamma_{rs}\big\|^{2}}_{\text{granular}},
\label{eq:variety_split}
\end{equation}
using $\mathbb{E}_{\Omega}[\omega_{rs}(\Omega)]=\gamma_{rs}$. The first term is the concentration of the expected network --- the allocation an infinite number of varieties would deliver, and the only one comparative advantage and trade costs act on. The second is the average distance between a drawn network and that limit, and is strictly positive whenever $N_{s}$ is finite.

Both terms have closed forms. Writing $\omega_{r'rs}(\Omega)=\sum_{j=1}^{N_{s}}v_{jrs}(\Omega)\,I_{jr's}(\Omega)$, with $v_{jrs}$ the share of buyer $r$'s sector-$s$ spending going to variety $j$ and $I_{jr's}=\mathbf 1\{r'\text{ wins }j\text{ for }r\}$,
\begin{equation}
\mathbb{E}_{\Omega}\big[H_{rs}(\Omega)\big]=\|\gamma_{rs}\|^{2}+\mathcal{V}_{s}\big(1-\|\gamma_{rs}\|^{2}\big),
\qquad
\mathcal{V}_{s}=\mathbb{E}_{\Omega}\Big[\textstyle\sum_{j}v_{jrs}^{2}(\Omega)\Big].
\label{eq:granular_buyer}
\end{equation}
The term $\mathcal{V}_{s}$ is the concentration of a buyer's sector-$s$ input bill across varieties, so that $1/\mathcal{V}_{s}$ is the effective number of varieties purchased; it depends only on $(\theta,\nu_{s},N_{s})$, not on the buyer's location.\footnote{Under Fr\'echet the price at which a variety is won is independent of the identity of its winner, $G_{r'rs}(p)=G_{rs}(p)$, so the $N_{s}$ champion prices are i.i.d.\ draws from a common distribution whose scale $\Phi_{rs}^{-1/\theta}$ alone carries $T$ and $\tau$. Expenditure shares $v_{\rho rs}=p_{\rho rs}^{1-\nu_{s}}/\sum_{j}p_{jrs}^{1-\nu_{s}}$ are homogeneous of degree zero in prices, so that scale cancels: geography shifts the level of the champion-price distribution, common to all of a buyer's varieties, and leaves its relative dispersion untouched. The same independence makes \eqref{eq:granular_buyer} exact rather than an approximation. Both results rely on firms sourcing on the expected price index rather than the realised CES aggregate, following \citet{lenoir2023search}. The restriction is testable rather than assumed: measured across buyers, the range of $\mathcal{V}_{rs}$ does not exceed $0.0052$ in motor vehicles and $0.0126$ in aerospace.} Geography and granularity therefore enter \eqref{eq:granular_buyer} through separate terms, and every counterfactual below moves the first alone.

The granular term has a simple reading. Granularity forces the system to produce exactly $N_{s}$ winners: each variety is won by one region, drawn with probability $\gamma_{r'rs}$, so a region's realised share is a weighted count of wins and the variance of that count is what the Herfindahl adds up. Under equal weights the variance is $\gamma_{r'rs}(1-\gamma_{r'rs})/N_{s}$; under CES the wins are weighted by expenditure and $\mathcal{V}_{s}$ replaces $1/N_{s}$, but the shape is unchanged.\footnote{\emph{Appendix :} the mechanism is clearest when varieties carry equal expenditure weights, $v_{jrs}=1/N_{s}$. Then $N_{s}\omega_{r'rs}(\Omega)$ is literally the number of varieties region $r'$ wins, which is $\operatorname{Binomial}(N_{s},\gamma_{r'rs})$, and since $\mathbb{E}[K^{2}]=N\gamma(1-\gamma)+N^{2}\gamma^{2}$ for $K\sim\operatorname{Bin}(N,\gamma)$, summing over $r'$ and using $\sum_{r'}\gamma_{r'rs}=1$ gives $\mathbb{E}_{\Omega}[H_{rs}(\Omega)]=H(\gamma_{rs})+N_{s}^{-1}[1-H(\gamma_{rs})]$. Under CES a cheaper winner takes a larger share, so $N_{s}\omega_{r'rs}$ is a \emph{weighted} sum of winner indicators rather than a count and that expression is not literally a binomial second moment; \eqref{eq:granular_buyer} is the general case, with $\mathcal{V}_{s}=1/N_{s}$ recovered under equal shares and $\mathcal{V}_{s}\ge 1/N_{s}$ always by Cauchy--Schwarz, so reading granularity as a count \emph{understates} it. The closed form for $\mathcal{V}_{s}$ has no content at this calibration and we measure it from the draws instead: the variety weights are Pareto with tail index $\kappa_{s}=\theta/(\nu_{s}-1)$ and $\mathcal{V}_{s}$ would be $\Xi(\kappa_{s})/N_{s}$ with $\Xi(\kappa)=\Gamma(1-2/\kappa)/\Gamma(1-1/\kappa)^{2}$, but $\theta=1$ and $\nu_{s}=1.5$ give $\kappa_{s}=2$ \emph{exactly}, the boundary at which $\Xi$ diverges, so the variance of the variety weights is not finite and a single variety can carry a non-vanishing share of a buyer's sector spending.} Each region's contribution vanishes at $\gamma_{r'rs}=0$ and $\gamma_{r'rs}=1$ and peaks in between, so granularity bears on the regions a buyer sources from occasionally, not on those it always or never uses. It is a large force when comparative advantage leaves the sourcing probabilities spread out, and none at all when comparative advantage has already settled the winner.
 
\paragraph{How much does granularity add?}
A great deal, and through one channel rather than two. At the estimates the median buyer's portfolio carries $\mathbb{E}_{\Omega}[H_{r}]=0.118$ in motor vehicles and $0.207$ in aerospace, against $H(\gamma_{r})=0.017$ and $0.055$ for the same buyer's expected network: with finitely many varieties the supplier portfolio is $6.9$ times as concentrated in motor vehicles and $3.7$ times as concentrated in aerospace as the allocation the model would deliver in the continuum. In effective numbers of supplier commuting zones, $58$ and $18$ become $8.5$ and $4.8$. Which of the two factors of $\mathcal{V}_{s}\,[1-H(\gamma_{rs})]$ produces that is settled by the fact that the second is bounded above by one. The slack $1-H(\gamma_{r})$ is $0.983$ and $0.945$ --- within $1.7$ and $5.5$ percent of its ceiling --- so comparative advantage has concentrated the expected network far too little to crowd chance out, and the granular term is very nearly $\mathcal{V}_{s}$ itself. The measured effective number of varieties, $1/\mathcal{V}_{s}=9.7$ and $6.4$, is therefore almost the whole of the realised statistic: $\mathcal{V}_{s}$ accounts for $87$ and $76$ percent of $\mathbb{E}_{\Omega}[H_{r}]$, and what remains for the geography is the last tenth or quarter. The reading is that with $\hat N_{s}$ between $12$ and $24$ and CES handing the cheapest winner a disproportionate share --- the effective count is about $0.6\,N_{s}$ --- a buyer buys a sector's inputs from of order ten distinct varieties, each won by one region, so its Herfindahl is close to a count of varieties wearing a region's label. That is a property of the object and not of the calibration: a Herfindahl weights regions by euros, and with finitely many varieties the euros pile onto whichever few varieties came out cheap. It is also the caution that governs everything below, including Figure~\ref{fig:buyer_concentration_region}: no intervention on $\hat T$ or $\alpha$ can move a statistic more than the tenth or quarter of it those forces own.

\paragraph{What the two forces do to it.}
Figure~\ref{fig:buyer_concentration_region} switches each force off buyer by buyer, as percentage deviations of $\mathbb{E}_{\Omega}[H_{r}]$ from the estimated economy. Three things read off it. First, \emph{comparative advantage is what concentrates the portfolio and trade costs do essentially nothing}: equalising $\hat T$ lowers concentration by $7.6$ percent in motor vehicles and $17.7$ in aerospace, while removing distance costs moves it by $+1.9$ and $-1.8$ percent --- not merely small but not even of a common sign across the two industries. Second, \emph{the two forces together do no more than comparative advantage alone}. The uniform benchmark, in which nothing selects among cells, sits at $-7.9$ and $-18.8$ percent, so removing $\hat T$ already accounts for $95$ and $94$ percent of the joint effect; the two are substitutes, not complements, in producing concentration. Third, and this is what the decomposition adds, \emph{those percentages badly understate what comparative advantage does to the underlying network}. On $H(\gamma_{r})$ alone, equalising $\hat T$ takes the median buyer from $0.017$ to $0.007$ and from $0.055$ to $0.018$ --- falls of $61$ and $68$ percent, essentially all the way to the unselective benchmark. The realised statistic moves by an eighth of that because the granular term absorbs the difference: as $H(\gamma_{rs})$ falls the slack rises toward one and $\mathcal{V}_{s}\,[1-H(\gamma_{rs})]$ rises with it, from $0.101$ to $0.102$ and from $0.148$ to $0.152$, offsetting most of the structural collapse. Chance fills in the dispersion that comparative advantage vacates, and the two are substitutes exactly as the previous paragraph's $1-H(\gamma_{rs})$ factor requires.\footnote{The $\alpha=0$ column should not be read as an economic effect. Setting $\alpha=0$ moves the structural term by $-0.0010$ and $+0.0005$ in levels --- distance is the weaker of the two forces here --- while the accompanying movement in $\mathcal{V}_{s}$, from $0.103$ to $0.106$ in motor vehicles and from $0.157$ to $0.154$ in aerospace, moves the granular term by $+0.0032$ and $-0.0029$. The second channel is three to six times the first and therefore sets the sign, and those $\mathcal{V}_{s}$ movements are of the order of the simulation error at $50$ replications, which is why the column changes sign between the industries. The figure's content is the $\hat T$ contrast and the benchmark.} One further feature of the figure is a result rather than a defect: the bars are nearly identical across buyers in motor vehicles, where the realised effective number ranges only from $8.4$ to $8.5$, and vary visibly in aerospace ($4.4$ to $5.0$). Granularity makes buyers look alike, because $\mathcal{V}_{s}$ is a sector object common to all of them and is most of what the statistic is made of, while the expected networks it overlays are far more dispersed across buyers.
 
\begin{figure}[H]
    \centering
\caption{Concentration of the supplier portfolio per buyer under each scenario\label{fig:buyer_concentration_region}}
    \begin{tabular}{cc}
    Motor vehicles&Aerospace\\
        \includegraphics[width = 0.5\linewidth]{figures/buyer_concentration_region_auto_mu2.pdf} &
        \includegraphics[width = 0.5\linewidth]{figures/buyer_concentration_region_aero_mu2.pdf}
    \end{tabular}
    \caption*{\footnotesize \emph{Notes}: $\mathbb{E}_{\Omega}[H_{r}(\Omega)]$ for each shocked commuting zone under each counterfactual regime, as a percentage deviation from \emph{Both forces}, the estimated allocation; the level it is taken against is printed beside each row, and the figure is read row for row against Figure~\ref{fig:counterfactual_distance_region}. The concentration itself rather than its reciprocal, because \eqref{eq:granular_buyer} is additive in $H$ and the distance between two regimes is then the change in the granular term; effective numbers are not additive. The infinite-variety limit is \emph{not} drawn: it sits at $-85$ and $-73$ percent, an order of magnitude beyond the two counterfactuals, and on the same axis would compress them onto the zero rule. That comparison is the previous paragraph, and it is the larger of the two.}
\end{figure}

\paragraph{How much of the shock does the region keep?}
How far the euro travels is not the same question as how much of it the shocked region retains, and it is the second that the empirical part of this paper measures. Let $$L_{r}(d)=\frac{\sum_{r':\,d_{r'r}\le d} X_{r'r}}{\sum_{r'} X_{r'r}}$$ be the share of the upstream response accruing within $d$ kilometres of the origin of the shock. We report it at $d=200$~km, the radius at which the reduced-form evidence of Section~\ref{sec:reduced_form_evidence} is measured; Appendix~\ref{app:radius} gives the whole profile $L_{r}(d)$ and confirms that the comparison is a level difference rather than a crossing. Figure~\ref{fig:counterfactual_local_share} reports it region by region under each regime and Figure~\ref{fig:local_share_map} maps it, since $L_{r}(d)$ is a number attached to a place.
 
Here the two industries separate most sharply, and the sign of the separation is the section's main result. In motor vehicles the two forces \emph{reinforce} each other: the local share falls from 23.8\% to 18.2\% when $\hat T$ is equalised and to 19.4\% when $\alpha=0$, so removing either disperses the shock. In aerospace they work \emph{against} each other. Switching off distance costs 6.7 percentage points (21.4\% to 14.7\%), so proximity is what keeps an aerospace shock local; but equalising $\hat T$ \emph{raises} the local share, from 21.4\% to 24.5\%. Comparative advantage in aerospace is not a localising force. It takes the response away from the region in which the shock originates and sends it to the hubs, so each aerospace region retains about three percentage points less of its own shock than it would in an economy with no differences in comparative advantage at all. That is the co-location term of Table~\ref{tab:variance_decomposition} arriving where it can be seen, on the euros a region keeps.
 
This is also why the two statistics are both needed, rather than one being a robustness check on the other. The mass taken off Toulouse and Île-de-France lands partly next to each buyer, which the local share records, and partly at an intermediate distance, which offsets it in the mean. How far the shock travels and who keeps it are separate questions, and in aerospace comparative advantage answers only the second.
 
The medians of $L_{r}(200)$ are indistinguishable --- 21.3\% in motor vehicles against 19.7\% in aerospace --- and the dispersion behind them is again what differs. In motor vehicles the local share ranges from 15.5\% to 43.7\%: shocks diffuse widely, but in much the same way wherever they start, and no region is left without a substantial local response. In aerospace the range runs from 3.1\% to 56.3\%. A shock to Toulouse is retained more than half within 200~km; a shock to most other aerospace regions is transmitted almost entirely elsewhere.
 
\begin{figure}[H]
    \centering
\caption{Local incidence with one force switched off\label{fig:counterfactual_local_share}}
    \begin{tabular}{cc}
    Motor vehicles&Aerospace\\
        \includegraphics[width = 0.5\linewidth]{figures/counterfactual_local_share200km_auto_mu2.pdf} &
        \includegraphics[width = 0.5\linewidth]{figures/counterfactual_local_share200km_aero_mu2.pdf}
    \end{tabular}
    \caption*{\footnotesize \emph{Notes}: $L_{r}(200)$, the share of the upstream response accruing within 200~km of the origin of the shock, one group of bars per shocked commuting zone, ordered by the realised local share. Regimes as in Figure~\ref{fig:counterfactual_distance_region}.}
\end{figure}
 
\begin{figure}[H]
    \centering
\caption{Where a downstream shock lands\label{fig:local_share_map}}
    \begin{tabular}{cc}
    Motor vehicles&Aerospace\\
        \includegraphics[width = 0.5\linewidth]{figures/local_share_map200km_auto_mu2.pdf} &
        \includegraphics[width = 0.5\linewidth]{figures/local_share_map200km_aero_mu2.pdf}
    \end{tabular}
    \caption*{\footnotesize \emph{Notes}: Each shocked downstream commuting zone is coloured by $L_{r}(200)$, the share of the upstream sales generated by a demand shock in that zone that accrues within 200~km of it. Commuting zones hosting no downstream producer are shaded in grey: they originate no shock. Both panels use the same colour scale.}
\end{figure}
 
\paragraph{Uncertainty.} The counterfactual quantities are reported as point estimates. They are functions of $\hat\alpha\hat\theta$ and $\hat T$ and inherit the sampling uncertainty of the estimation.\swann{A parametric bootstrap over $(\hat\alpha,\hat T)$, redrawing from the estimated sampling distribution and recomputing each counterfactual, is the standard way to attach an interval to these numbers and is not yet run.}
 
\paragraph{Taking stock.} This is the efficiency--stability trade-off of the introduction made quantitative, and the two industries sit on opposite sides of it. The concentration of high-productivity aerospace suppliers around Toulouse and Paris raises the competitiveness of the downstream industry, but it also means that a shock originating anywhere in the industry converges on the same two areas: an aerospace buyer's supplier portfolio spans an effective 18 commuting zones against 58 in motor vehicles --- 28\% and 38\% of what an unselective allocation would give them --- and 95 to 99 per cent of where the euro lands is common to every buyer, so those supplier regions hold undiversified clienteles and the shocks hitting different buyers pile up on them rather than offsetting each other. Finite varieties sharpen both margins rather than blurring them: in a single year the portfolio collapses to an effective 4.9 zones in aerospace and 8.2 in motor vehicles, and the chance winners are common across buyers in eight to nine cases out of ten, against a floor of one in twenty. The shock does not travel further than it otherwise would --- the hub sits at about the distance the euro would cover anyway, so $d_{r}$ barely registers the concentration --- it travels to the same place, whoever is hit, and the regions in which it originates retain three percentage points less of it than they would in an economy with no differences in comparative advantage at all. Motor vehicles, whose upstream comparative advantage is co-located with a more dispersed set of downstream producers, buys a 65~km shorter supply chain with that co-location and leaves no region without a local response.
 

\subsection{Spatial consequences of military procurements and EV subsidies}

\paragraph{Military Procurement in the Aerospace Sector.} The rise in geopolitical tensions has led several countries, including France, to increase their defense spending. As a consequence, firms in the defense sector are exposed to positive demand shocks that are largely largely driven by geopolitical considerations rather than firm-level conditions. We next use such shocks to evaluate the model's predictions on the spatial diffusion of demand-driven changes in activity.

Under the 2019–2025 \textit{Loi de programmation militaire} (LPM), the French government commissioned 169 H160M helicopters from Airbus Helicopters. The contract covers design, production, and maintenance and is estimated at approximately €10 billion over ten years. In the context of our model, this policy constitutes a large, plausibly exogenous demand shock to a single dwonstream producer, which propagates through the production network to its suppliers.

This shock is particularly informative for our analysis because of its spatial dimension: Airbus Helicopters’ headquarters are located in Marignane, north-west of Marseille. As documented in Section \ref{sec:data_stylized_facts}, this region is not among the most specialized in aerospace equipment, implying that the initial impact is geographically concentrated but that equilibrium effects are likely to diffuse across space through supplier linkages. In addition, the contract is sizeable, representing approximately twice the firm’s 2019 annual revenue, which ensures that spillover effects through the production network are economically meaningful.
\swann{Ajouter une footnote pour justifier que l'on peut rester en PE étant donné la taille du choc. } \isa{More generally, we have not mentioned our assumptions on input costs when discussing the model's estimation. We need to add something there, before eventually mentioning GE effects of procurement here.}

Within our model, the contract generates €470 million in annualized incremental sales for direct suppliers of Airbus-Helicopters. \isa{Not sure about this statement. You mean that 10 billion over 10 years is approximately 1 billion a year, which generates an addition 470 million per year in upstream sectors? } Due to the geographic isolation of Marseille--Aubagne, the local economy relies heavily on supply chains centered in Toulouse and Paris. This structural dependence ensures that the shock propagates toward distant spatial nodes, with 89\% of the incremental sales being captured by suppliers in those two regions, leaving a mere 1\% for firms within the Marseille--Aubagne area. \isa{how can we be slightly more precise here? We already know from Figure 6 that the amplification is roughly .5, ie 500 million per year and that less than 2\% of this goes to Marseille-Aubagne. So maybe we can discuss more which exact CZ benefit? I understand that this will be Toulouse and Ile de france but how much for each? Other CZ? maybe a map here?}

%Consequently, the two-hub structure of the aerospace industry generates strong spatial comovement across regions. Shocks to downstream activity originating in Toulouse and Île-de-France are largely absorbed locally, leading to a concentration and amplification of fluctuations within these hubs. Conversely, shocks arising in peripheral regions are transmitted through production linkages toward these central nodes, rather than remaining locally contained. This asymmetric propagation implies that both local and remote shocks contribute to fluctuations in hub regions, thereby synchronizing economic activity across space and reinforcing aggregate comovement.



\paragraph{EV subsidies in the Car industry.}


For the motor vehicles industry, we quantify the spatial consequences of the 2024 reform of the French EV environmental score. Prior to 2024, all battery-electric vehicles meeting price and weight thresholds qualified for a purchase subsidy of approximately \euro{5,000}. In January~2024, France introduced a life-cycle emissions threshold---the \textit{score environnemental}---that disqualified models whose manufacturing-stage carbon footprint exceeded a fixed ceiling. Because upstream emissions correlate strongly with the location of assembly and battery production, the reform effectively excluded most Chinese-assembled models. Roughly 30\% of 2023 EV sales lost eligibility overnight, corresponding to a price increase of approximately \euro{4,260} per disqualified vehicle.
 
Using exhaustive French registration data matched to assembly plants and battery origins, \cite{MalgouyresMayerNolden2025} estimate a three-level nested logit demand system and show that the reform triggered a sharp reallocation of sales toward European-assembled vehicles, while reducing overall EV penetration by approximately 0.9~percentage points. Their difference-in-differences estimates imply that losing the bonus reduces sales of affected models by roughly 60\%, with an own-price semi-elasticity for EVs of $\hat{\beta}_{EV} = -0.154$ per thousand euros. The nested logit parameters---estimated by combining the quasi-experimental price variation with US second-choice survey moments from \cite{AllcottEtAl2024}---are $\hat{\alpha} = 0.087$, $\hat{\mu}_3 = 0.648$, and $\hat{\mu}_{4,E} = 0.569$.
 
We rely on these estimates to construct region-specific demand shocks affecting downstream motor vehicle producers, and feed them through our structural model to quantify the upstream spatial diffusion. The shock exploits regional heterogeneity in the production composition across three vehicle categories---eligible EVs, non-eligible EVs, and ICE vehicles---so that regions specializing in the assembly of disqualified models experience larger negative demand shifts. A detailed description of the procedure is provided in Appendix~\ref{app:ev_ecoscore}.
 
Overall, the estimated size of the shock amounts to X euros, and the additional upstream sales implied by our model are equal to X euros.


\begin{comment}
    
\paragraph{Taking stock} We have shown that the spatial organization of production networks concentrates the transmission of shocks toward a small number of hub regions. In the aerospace industry, a large share of suppliers serving multiple downstream regions are located in Toulouse and Île-de-France, implying that shocks originating in other regions are disproportionately propagated toward these hubs. In this sense, hub regions act as aggregators of network-driven fluctuations, as their activity loads on shocks affecting a broad set of downstream locations.

Our model focuses on the motor vehicle or aerospace production network and therefore captures this exposure within a given industry. Combining this mechanism with the empirical evidence on the sectoral composition of regions documented in Section~\ref{sec:data_stylized_facts} provides additional insight into the implications for local volatility. In particular, Toulouse is highly specialized in aerospace activity, whereas Île-de-France is significantly more diversified across sectors. As a result, the network-driven exposure identified in the model is likely to translate into a higher sensitivity of local economic activity to aerospace shocks in Toulouse than in Île-de-France.

This observation raises a policy-relevant question: to what extent can the spatial structure of production networks be reshaped in order to reduce the exposure of key regions to sector-specific shocks, while preserving the benefits of industrial concentration? In the next section, we use the model to address this question by analyzing how changes in trade costs and regional comparative advantage affect the organization of production networks and the resulting exposure of regions to downstream shocks.

 


\subsubsection{Reorganizing production network}


\swann{Ci dessous je donne le détail des calculs mais on pourrait en renvoyer une partie en Appendix.}

\subsubsection*{Theoretical results}

The model implied production network is the outcome of trade costs and sector - region comparative advantage. A policymaker can influence the production network structure either by reducing trade cost between region (investment in infrastructure) of increasing the comparative advantage of regions (through cluster-based policies). We seek in this section to understand how the reduction of trade costs and the increase of comparative advantage is going to i) reorganize the network, ii) affect the comovement between regions. We first expose theoretical results and then run simulations to quantify those effects. 

A region is composed by downstream firms and upstream firms. Sales change of upstream firms is totally determined by the transmission of downstream shocks: $d\ln x_{it} = \sum_r a_{ri}^D \delta_{r,t}^D$. Consequently, the spatial comovement between two regions is given by: 

$$\Cov (r',l) = M_{r'.}\Sigma_{dd}M_{l.}^T$$

with $M_{r'r} = \sum_{s} w_s^{r'} w_{d}^{sr'}w_{r}^{dsr'}$ the network matrix and $\Sigma_{dd}$ the variance covariance matrix of downstream shocks. $w_s^{r'}$ is the weight of $s$ in $r'$, $w_{d}^{sr'}$ is the share of sales of $s$ in $r'$ to the downstream industry and $w_{r}^{dsr'}$ is the share of those sales toward downstream region $r$, the outcome of the estimated model:


\begin{equation}
w_r^{sr'd} \;=\; \frac{\gamma_{r'rs}\,X_{rs}}{\sum_{r''}\gamma_{r'r''s}\,X_{r''s}},
\end{equation}

with $X_{rs}$ the purchase of inputs from sector $s$ by downstream firm in $r$. Any change in trade cost or comparative advantage will affect upstream sales portfolio and shift  $w_{r}^{dsr'}$, affecting the spatial comovement $\Cov(r',l)$. To assess such effect, we first derive $\gamma_{r'rs}$ with respect to $\tau_{krs}$ and $T_{ks}$ for any upstream region $k$:

\begin{align}
    \frac{d\ln \gamma_{r'rs}}{d\ln \tau_{krs}} &= -\theta (\mathds{1}\{r'=k\}- \gamma_{krs}) \\
    \frac{d\ln \gamma_{r'rs}}{d\ln T_{ks}} &=  (\mathds{1}\{r'=k\} - \gamma_{krs}) 
\end{align}

A decrease of trade cost or increase in comparative advantage has two distinct effects. If $k = r'$ it will increase the purchase from this $k$, the less if $r$ is already purchasing from $k$. If $k \neq r'$, it will decrease the purchase since $r'$ become relatively less competitive. Turning to $w_r^{dsr'}$ we have for any variable $y$: 

\begin{align}
    \frac{d\ln w_r^{dsr'} }{d\ln y} = \frac{d\ln \gamma_{r'rs}}{d\ln y} -  \sum_{r^{''}} w_{r^{''}}^{dr's}\frac{d\ln \gamma_{r'r^{''}s}}{d\ln y} + \frac{d\ln X_{rs}}{d\ln y} -  \sum_{r^{''}} w_{r^{''}}^{dr's}\frac{d\ln X_{r^{''}s}}{d\ln y}
\end{align}


For the remaining we consider $X_{rs}$ constant for all $r$. \swann{Je fais cette hypothèse car ça permet de simplifier les calculs dans un premier temps. Je pense que c'est important dans le modèle final de ne pas la faire. }Since for all $r^{''} \neq r$ we have $\frac{d\ln \gamma_{r'r^{''}s}}{d\ln \tau_{krs}} = 0$:

\begin{align}
\frac{d \ln w_r^{sr'd}}{d \ln \tau_{krs}} &\;=\; \bigl(1 - w_r^{sr'd}\bigr)\,\frac{d \ln \gamma_{r'rs}}{d \ln \tau_{krs}} \;=\; -\theta\bigl(\mathds{1}\{r'=k\}-\gamma_{krs}\bigr)\bigl(1-w_r^{sr'd}\bigr),
\label{eq:w_elast_tau_destr}\\[4pt]
\frac{d \ln w_{r''}^{sr'd}}{d \ln \tau_{krs}} &\;=\; -w_r^{sr'd}\,\frac{d \ln \gamma_{r'rs}}{d \ln \tau_{krs}} \;=\; +\theta\bigl(\mathds{1}\{r'=k\}-\gamma_{krs}\bigr)\,w_r^{sr'd}, \qquad r''\neq r.
\label{eq:w_elast_tau_otherdest}
\end{align}


Reducing trade cost between $k$ and $r$ is going to concentrate the portfolio around $r$ if $k = r'$ and away from $r$ otherwise. For the comparative advantage the effect is a bit more complicated. Because $T_{ks}$ affects $\gamma_{r'r''s}$ for all destinations $r''$ (not just one), the chain rule yields a portfolio-wide reallocation:
\begin{align}
\frac{d \ln w_r^{sr'd}}{d \ln T_{ks}} 
&\;=\; \frac{d \ln \gamma_{r'rs}}{d \ln T_{ks}} - \sum_{r''}w_{r''}^{sr'd}\frac{d \ln \gamma_{r'r''s}}{d \ln T_{ks}} \notag\\
&\;=\; \bigl(\mathds{1}\{r'=k\}-\gamma_{krs}\bigr) - \sum_{r''}w_{r''}^{sr'd}\bigl(\mathds{1}\{r'=k\}-\gamma_{kr''s}\bigr) \notag\\
&\;=\; \sum_{r''}w_{r''}^{sr'd}\,\gamma_{kr''s} - \gamma_{krs}
\label{eq:w_elast_T}
\end{align}
The indicator terms cancel in the third line: the result is independent of whether $r' = k$. This is sharper than previous elasticity, which depends on the relationship between origin $r'$ and the focal $k$.

\medskip\noindent
\textit{Mechanism.} A rise in $T_{ks}$ produces a \emph{uniform within-portfolio reallocation} for every origin $r'$: allocation shifts toward destinations $r$ where $k$ is currently under-represented (small $\gamma_{krs}$ relative to the portfolio-average $\sum_{r''}w_{r''}^{sr'd}\gamma_{kr''s}$), and away from destinations where $k$ is over-represented. This means cluster policy in $k$ has a \emph{symmetric effect across all origins} — every supplier network reorganizes toward the same destinations regardless of where the supplier is located. This is in sharp contrast to the trade cost reduction, where infrastructure linking $k$ to a specific destination $r$ produces origin-specific responses.\\


\swann{La formule (13) n'est pas intuitive au début. Pour bien la comprendre, on peut placer trois points sur une ligne. Le point de gauche et celui du milieu sont des régions avec une upstream et downstream firm. Le point de droite a uniquement une downstream firm. On augmente $T_{sr}$ au milieu. Toutes les downstream vont + sourcer du milieu mais celles aux extrémités relativement moins que celle du milieu. La redirection du flux va concentrer le portfolio du point de gauche sur les downstream des extrêmes. On peut ajouter une footnote pour l'expliquer.}

Turning to spatial comovements, abstracting from any across sector spillover we have for any variable $y$: 


\begin{equation}\label{eq:V_elast_tau}
    \frac{d \Cov(r',l)}{d \ln y} \;=\; \sum_{r}\Bigl[\bigl(\Sigma M_{l\cdot}^\top\bigr)_r M_{r'r}\,\frac{d \ln \ln M_{r'r}}{d \ln y}                 + \bigl(\Sigma M_{r'\cdot}^\top\bigr)_r M_{lr}\,\frac{d \ln M_{lr}}{d \ln y}\Bigr].
\end{equation}

When $y$ is $\tau_{krs}$, a decrease in $\tau_{krs}$ lead to an increase in $M_{kr}$ and a decrease in $M_{r'r}$ for $r'\neq q$. Consequently, the total effect on spatial comovement is ambiguous. If the reorganization of the network is such that $r'$ and $l$ are more exposed to similar shocks then it increases the spatial comovement between them. If downstream shocks are iid, $\Cov(k,r)$ increases and $\Cov(k,l)$ decreases for $l \neq r$. Reversely,  $\Cov(r',r)$ decreases for $r' \neq r$ and $\Cov(r',l)$ increases for $l \neq r$. \\

When $y$ is $T_{ks}$, the total effect is ambiguous and depends on portfolio reorganization interacted with the structure of the covariance matrix. 


\subsubsection*{Optimal cluster-based policy: a portfolio result}

We use the elasticities above to derive the optimal allocation of a fixed budget $B$ for cluster-based industrial policy.

\swann{Cet exercice est pour moi intéressant conventuellement, car il permet de reboucler avec la motivation initiale et de réfléchir à la manière de construire spatialement les chaînes de valeurs au sein d'un pays (Résilience des systèmes de production). Par exemple, compte tenu de la distribution des usines de l'automobile, est ce que c'est optimal de faire la vallée de la batterie dans les Hauts-de-France ?  Mais si jamais le coût de la volatilité est faible (et c'est ce que je comprends de la littérature) alors il n'est pas trop intéressant quantitativement parlant. }

\swann{Ici le planner internalise le coût de la volatilité. L'entreprise ne l'internalise pas quand elle forme ses liens. Si jamais on avait fait un modèle multi pays je pense que cette hypothèse aurait été problèmatique mais comme on est within country je pense que c'est ok. }

\paragraph{Setup.}
A planner allocates a budget $B$ across regions to subsidize comparative advantage. Let $b_{r}\geq 0$ be the budget allocated to region $r$, with $\sum_r b_r \leq B$. The induced productivity satisfies $T_{rs}(b_r) = \widehat T_{rs}\cdot f(b_r)$, where $\widehat T_{rs}$ is the estimated baseline and $f$ is a strictly increasing, strictly concave function with $f(0)=1$ and $f'(0) > 0$. Concavity captures the assumption of diminishing returns to local R\&D and infrastructure investment. \swann{Je n'ai pas de référence en tête mais j'imagine que c'est standard.}

The planner's objective combines an efficiency term (the negative of an aggregate downstream price index, weighted by destination expenditure) and a stability constraint (an upper bound on aggregate variance). Formally: 



\begin{equation}
\min_{\{b_r\}_{r=1}^R}\;\; P\bigl(\{T_{rs}(b_r)\}\bigr) \quad\text{s.t.}\quad \sum_r b_r \leq B,\;\; \Omega\Sigma\Omega^T \leq \bar V,\;\; b_r \geq 0,
\label{eq:planner_problem}
\end{equation}
where $P$ is the downstream industry price index, $\Sigma$ the region level variance covariance matrix, $\Omega$ planners weights on those comovements and $\Bar{V}$ the constraint on the aggregate variance. Investing in a specific location increases $T_{sr}$ and therefore lowers the marginal cost of production of downstream firms sourcing from this region. Consequently, the aggregate price index is reduced. However, portfolio reallocation might increase the spatial comovement affecting local and aggregate stability.




\paragraph{The price-index elasticity.}
In Eaton--Kortum, the sectoral price index satisfies $p_{rs}^{-\theta}\propto \Phi_{rs}$ (up to constants from the gamma function). Differentiating:
\begin{equation}
\frac{\partial \ln p_{rs}}{\partial \ln T_{ks}}  \;=\; -\frac{1}{\theta}\frac{\partial \ln \Phi_{rs}}{\partial \ln T_{ks}} \;=\; -\frac{\gamma_{krs}}{\theta}.
\label{eq:price_elast_T}
\end{equation}

Aggregating across sectors and destinations, the welfare-relevant price elasticity is
\begin{equation}
\frac{\partial \ln P}{\partial \ln T_{ks}} \;=\; -\frac{1}{\theta}\sum_{r}\omega_r^{\mathcal{P}}\,\Omega^s\,\gamma_{krs} \;\equiv\; -\frac{1}{\theta}\,\bar\gamma_{ks},
\label{eq:welfare_elast_T}
\end{equation}
where $\omega_r^{\mathcal{P}}$ is destination $r$'s weight in welfare and $\bar\gamma_{ks} \equiv \sum_r \omega_r^{\mathcal{P}}\Omega^s\gamma_{krs}$ summarizes $k$'s \emph{network centrality} in supplying sector $s$. Centrally located regions with high average sourcing shares deliver higher marginal welfare returns per unit of comparative advantage.

\paragraph{ The variance elasticity.}
Define the variance elasticity $\mathcal{V}'_{ks} \equiv \partial \mathcal{V}/\partial \ln T_{ks}$, computed via \eqref{eq:V_elast_tau} applied to the appropriate variance aggregate. From §B.4, $\mathcal{V}'_{ks}$ is sign-ambiguous: negative when $k$ is a \emph{stabilizing} hub (its shock is uncorrelated or counter-correlated with the existing exposures of regions reallocating toward it), positive when $k$ is a \emph{concentrating} hub.

\paragraph{ First-order conditions.}
Let $\lambda \geq 0$ be the budget multiplier and $\mu \geq 0$ the multiplier on the variance constraint. The Lagrangian first-order condition for $b_r > 0$ is

\begin{equation}
\boxed{\;f'(b_r)\cdot\underbrace{\Bigl[\,\underbrace{\tfrac{1}{\theta}\bar\gamma_{rs}}_{\text{efficiency gain}}                  \;-\;
\underbrace{\mu\,\mathcal{V}'_{rs}}_{\text{variance cost}}\,\Bigr]}_{\text{Marginal value of subsidy in $r$}}
\;=\;\lambda
\quad\text{for all $r$ with $b_r > 0$.}
\;}
\label{eq:foc}
\end{equation}

The marginal value per euro of subsidy in region $r$ is the sum of two components: an efficiency gain proportional to $r$'s network centrality $\bar\gamma_{rs}$, and a variance cost proportional to $r$'s contribution to aggregate variance $\mathcal{V}'_{rs}$. Diminishing returns $f'(b_r)$ ensure that the optimal allocation spreads investment across multiple regions rather than concentrating in one.

\paragraph{ The interior characterization.}
Equation \eqref{eq:foc} delivers the central theoretical result of this section. Two regions $r$ and $r'$ receive interior allocations ($b_r, b_{r'} > 0$) if and only if their marginal values are equalized:
\begin{equation}
\frac{f'(b_r)}{f'(b_{r'})}
\;=\;
\frac{\frac{1}{\theta}\bar\gamma_{r's} - \mu\mathcal{V}'_{r's}}
     {\frac{1}{\theta}\bar\gamma_{rs} - \mu\mathcal{V}'_{rs}}.
\label{eq:interior}
\end{equation}
Concavity of $f$ ($f' > 0$, $f'' < 0$) implies that regions with higher marginal value receive larger allocations. The composition of "marginal value" is the substantive content of the result.

\begin{proposition}[Optimal cluster policy]
\label{prop:cluster_policy}
Under \eqref{eq:planner_problem} with strictly concave $f$, the optimal allocation $\{b_r^\ast\}$ favors regions $r$ that are jointly:

\begin{enumerate}
\item \emph{Centrally located in the production network of sector $s$} ($ \bar\gamma_{rs}$ large) — yielding high efficiency gains; and 
\item \emph{Diversifying} relative to the existing covariance structure of  ownstream shocks ($\mathcal{V}'_{rs} < 0$) — yielding stability gains.
\end{enumerate}

Conversely, regions that are already-saturated hubs of incumbent sourcing ($\gamma_{r\cdot s}$ near 1 in many destinations) deliver low marginal returns on both dimensions: efficiency returns are small because their share is bounded above, and the diversification benefit is negative because reinforcing them concentrates aggregate exposure.
\end{proposition}

\medskip\noindent
\textit{Mechanism and policy interpretation.} The optimal cluster-based policy is not the one that reinforces existing hubs. It is the one that develops regions that are \emph{moderately central but currently under-supplied}. These regions deliver two complementary benefits. First, by virtue of central location, an increase in their comparative advantage propagates to many downstream firms, lowering the price index across the network. Second, because their downstream-shock exposure is not yet large, drawing more sourcing toward them diversifies the existing portfolio of upstream regions, smoothing local business-cycle fluctuations. Reinforcing already-saturated hubs (e.g., Toulouse in aerospace) yields diminishing efficiency returns and adds variance through concentration. Reinforcing isolated peripheral regions (e.g., a CZ with no pre-existing linkages) yields high variance reduction in principle but low efficiency because the productivity gain is poorly diffused across the network. The optimum lies at intermediate centrality.

\medskip\noindent

\textit{Empirical implementation.} For each sector $s$, we compute $\bar\gamma_{rs}$ and $\mathcal{V}'_{rs}$ from the estimated model and report:
\begin{enumerate}
\item the cross-section of marginal welfare returns $\bar\gamma_{rs}$ across upstream regions;
\item the cross-section of variance returns $\mathcal{V}'_{rs}$ across the same regions;
\item the resulting optimal allocation $b_r^\ast$ under a calibrated $f(\cdot)$, for two specifications of the variance constraint ($\mathcal{V}$ = aggregate, $\mathcal{V}$ = local for selected regions).
\end{enumerate}
The headline statement is the comparison between $b_r^\ast$ and the allocation that would maximize efficiency alone ($\mu = 0$), which reinforces existing hubs. The variance-constrained optimum redirects a quantifiable share of the budget toward moderately-central, diversifying
regions.


\end{comment}